%======================================================================
\section{Introducción}\label{sec:intro}
%======================================================================

\subsection{Para quién es este documento y cómo leerlo}
Este documento está pensado para alguien que llega al proyecto sin conocerlo: un estudiante
de ingeniería o de matemática aplicada que sabe cálculo, álgebra lineal y ecuaciones
diferenciales ordinarias, pero que no necesariamente conoce la física de reactores, el
control predictivo ni la optimización numérica. Por eso cada concepto se introduce antes de
usarse, con un ejemplo cuando ayuda, y cada resultado se acompaña de la pregunta que
responde.

El texto tiene cuatro partes:
\begin{enumerate}[label=\textbf{Parte \Roman*.}, leftmargin=*]
  \item \textbf{El problema físico} (secciones~\ref{sec:reactor}--\ref{sec:accidente}): qué
        es el reactor SMDFR, cómo se describe su dinámica (neutrones y calor) y qué ocurre en
        el accidente que se quiere controlar.
  \item \textbf{El controlador} (secciones~\ref{sec:mpc}--\ref{sec:actual}): qué es un
        control predictivo, qué problema de optimización resuelve en cada instante y cómo lo
        resuelve hoy el método del proyecto.
  \item \textbf{Las herramientas} (secciones~\ref{sec:optimizacion}--\ref{sec:certteoria}):
        los fundamentos de optimización, integración numérica, linealización, incertidumbre y
        certificación de optimalidad, y las formulaciones que propone este trabajo.
  \item \textbf{Los resultados} (secciones~\ref{sec:replica}--\ref{sec:conclusiones}):
        qué se obtuvo con cada formulación, cuál es mejor, con qué garantía y qué falta.
\end{enumerate}
Quien ya conozca la física puede saltar a la parte~II; quien conozca el control predictivo,
a la parte~III. Los recuadros sombreados resumen la idea principal de cada sección. Al final
hay una tabla de notación, un glosario y las instrucciones para reproducir cada número.

\subsection{La problemática en breve}
El SMDFR (\emph{Small Modular Dual Fluid Reactor}) es un reactor de 100~MWt con combustible
de sal fundida y refrigerante de plomo líquido. Durante un accidente con pérdida de
refrigerante (LOCA), el inventario de plomo baja, el caudal del circuito primario cae y sube
la temperatura pico de la vaina de los tubos de combustible (PCT, \emph{peak cladding
temperature}). Un sistema auxiliar de enfriamiento (ACS) puede reponer refrigerante, pero
hay que decidir \emph{cuánto} y \emph{cuándo}: con poco, la PCT se acerca al límite de
seguridad; con mucho, se gasta un recurso limitado y se fuerza un componente que puede
fallar.

El proyecto aborda esa decisión con un \textbf{gemelo digital informado por riesgo}. Un
control predictivo (MPC) elige, cada 0.5~s, el caudal del ACS para los próximos 12.5~s.
Para predecir usa tres piezas: un modelo reducido del reactor, un pronóstico de la fuga
hecho por una red neuronal, y un árbol de eventos con las fallas y reparaciones posibles del
ACS. La seguridad se mide con un margen probabilístico (PSM).

El problema de optimización del MPC se resolvía con \textbf{metaheurísticas}
(Luus--Jaakola y algoritmo genético), porque cada evaluación del costo era una simulación de
caja negra en Simulink. Eso tiene tres consecuencias:
\begin{itemize}
  \item cada consulta tarda unos 800~s, cuando la meta del lazo es menos de 120~s;
  \item la acción aplicada depende de la semilla aleatoria del optimizador;
  \item no se sabe cuán lejos está la solución de la mejor posible.
\end{itemize}
Además, sin un problema tratable no se puede incorporar el riesgo dentro de la decisión,
porque eso multiplica el número de escenarios que hay que evaluar. Este documento propone
\emph{formulaciones monolíticas} del mismo problema: el modelo deja de ser una caja negra y
pasa a formar parte del problema de optimización como un conjunto de ecuaciones derivables.
Luego evalúa esas formulaciones y certifica la calidad de sus soluciones.

\subsection{Preguntas que responde este documento}
\begin{enumerate}[label=\textbf{Q\arabic*.}]
  \item ¿Se puede resolver el problema del MPC con exactitud y en tiempo real?
  \item ¿La solución obtenida es el óptimo global? ¿Con qué garantía?
  \item ¿Cuál método produce la mejor decisión, y por cuánto?
  \item ¿La incertidumbre del pronóstico de fuga cambia la decisión?
  \item ¿Las cotas del actuador tienen base física?
\end{enumerate}
Cada sección de resultados termina con un recuadro que responde una de ellas, y la
sección~\ref{sec:conclusiones} las reúne.

\subsection{Ruta de trabajo}
La idea central es escribir las ecuaciones del modelo reducido \emph{una sola vez}, en
forma algebraica y derivable. De esa fuente única salen el problema exacto, sus
linealizaciones y las cotas que certifican la calidad de las soluciones. El trabajo se
organizó en fases (tabla~\ref{tab:fases} y figura~\ref{fig:fases}).

\begin{table}[H]\centering\small
\begin{tabularx}{\linewidth}{@{}l Y Y l@{}}\toprule
Fase & Objetivo & Criterio de cierre & Estado\\\midrule
F0 Ecuaciones & Escribir el ROM término a término desde el archivo de Simulink & Cada término trazado a un bloque & Cerrada\\
F1 Réplica & ROM algebraico en Julia, con diferenciación automática & Error $\ll1$~K frente a Simulink & Cerrada\\
F2 NLP (N3) & P0 como programa no lineal monolítico & Comparado con LJ y GA en Simulink & 2 consultas\\
F3 QP (N1, N2) & P0 linealizado como QP; función de transferencia; certificado dual & Brecha frente a N3 medida & 2 consultas\\
F3c Certificación & Brecha de optimalidad global del P0 no lineal; comparación de decisiones & Brecha certificada pequeña & Pasos 1--4b\\
F4 Riesgo & Restricción probabilística, SAA, cuantiles & Una decisión que cambie por riesgo & Abierta\\
F5 Lazo cerrado & Las 61 consultas, multietapa & $<120$~s por consulta & Pendiente\\\bottomrule
\end{tabularx}
\caption{Fases del trabajo (detalle en \code{docs/monolithic\_formulation\_plan.md}).}
\label{tab:fases}
\end{table}

\begin{figure}[H]\centering
\resizebox{\linewidth}{!}{%
\begin{tikzpicture}[
  box/.style={draw=accent, line width=0.6pt, rounded corners=2pt, fill=soft,
              align=center, font=\small, minimum height=13mm, text width=24mm, inner sep=2pt},
  ref/.style={draw=black!50, dashed, rounded corners=2pt, align=center, font=\small,
              minimum height=10mm, text width=30mm},
  arr/.style={-{Stealth[length=2mm]}, line width=0.6pt, draw=ink},
  val/.style={-{Stealth[length=2mm]}, line width=0.5pt, draw=black!50, dashed},
  lab/.style={font=\footnotesize, text=black!70, align=center, inner sep=1.5pt}]
\node[box] (f0) at (0,0)      {\textbf{F0}\\Ecuaciones\\del \code{.slx}};
\node[box] (f1) at (3.6,0)    {\textbf{F1}\\Réplica\\algebraica};
\node[box] (f2) at (8.8,1.6)  {\textbf{F2}\\NLP no lineal\\(N3)};
\node[box] (f3) at (8.8,-1.6) {\textbf{F3}\\Linealización\\(N1, N2)};
\node[box] (f4) at (13.4,0)   {\textbf{F4}\\Riesgo\\(probab., SAA)};
\node[box] (f5) at (17.0,0)   {\textbf{F5}\\Lazo cerrado\\(61 consultas)};
\node[box] (f3c) at (13.4,-3.4) {\textbf{F3c}\\Certificación\\y comparación};
\node[ref] (sim) at (3.6,-3.6) {Simulink (ROM)\\referencia exacta};
\draw[arr] (f0) -- (f1);
\draw[arr] (f1.east) -- (5.0,0) |- node[lab, pos=0.75, above]{restricciones} (f2.west);
\draw[arr] (f1.east) -- (5.0,0) |- node[lab, pos=0.75, below]{jacobianos} (f3.west);
\draw[arr] (f2) -- node[lab, right]{óptimo de\\referencia} (f3);
\draw[arr] (f2.east) -- (11.2,1.6) |- (f4.west);
\draw[arr] (f3.east) -- (11.2,-1.6) |- (f4.west);
\draw[arr] (11.2,-1.6) |- (f3c.west);
\draw[arr] (f3c) -- node[lab, right]{garantía} (f4);
\draw[arr] (f4) -- (f5);
\draw[val] (sim.north) -- node[lab, right]{valida} (f1.south);
\draw[val] (sim.east) -| node[lab, pos=0.3, below]{re-simula cada acción} (f3.south);
\end{tikzpicture}}
\caption{Cómo se encadenan las fases. F0 y F1 producen un modelo derivable y validado. F2
lo usa como restricción (el problema exacto) y F3 como fuente de derivadas (el problema
lineal). F3c certifica la calidad de las soluciones. F4 y F5 construyen sobre un problema
que ya es tratable y verificado. Las flechas discontinuas son validaciones contra
Simulink.}
\label{fig:fases}
\end{figure}

\takeaway{\textbf{Idea del documento.} El controlador del reactor necesita resolver, cada
0.5~s, un problema de optimización sobre un modelo físico. Si el modelo se trata como caja
negra, sólo se puede buscar a ciegas. Si se escribe como ecuaciones, se pueden usar
derivadas: el problema se resuelve unas $10^4$ veces más rápido, de forma reproducible y con
garantías sobre la calidad de la solución.}
